

\chapter{Python and pandas}\label{ch:iv3:python-and-pandas}

\begin{ivtrains}
Knowing what the interpreter really does, then what a library does to your data. The stems say which version they mean where it matters: here CPython 3.10, NumPy 2.2 and pandas 2.3. The pandas questions are mostly about getting the right number, since a fast wrong answer is the usual failure.
\end{ivtrains}

\begin{question}[$\star$ \iqroles{QD, Dev, QR, DS} \iqfirm{any} \iqlevel{junior}]\label{iq:iv3:python-and-pandas:1}
What four lines does \cref{lst:iv3:pp:default} print on CPython, and why? Rewrite \texttt{add} so that each call without a list starts from an empty one.
\omcode{code/interviews-3/02-python-and-pandas/python/pp_default.py}{2}{10}{A default argument.}\label{lst:iv3:pp:default}
\end{question}

\begin{question}[$\star$ \iqroles{QD, Dev, QR, DS} \iqfirm{any} \iqlevel{junior}]\label{iq:iv3:python-and-pandas:2}
What two lines does \cref{lst:iv3:pp:alias} print? Fix each so that every key and every row has its own list. Why is \texttt{dict.fromkeys(names, 0)} safe, when \texttt{dict.fromkeys(names, [])} is not?
\omcode{code/interviews-3/02-python-and-pandas/python/pp_alias.py}{1}{7}{Two ways to share a list.}\label{lst:iv3:pp:alias}
\end{question}

\begin{question}[$\star$ \iqroles{QD, Dev, DS} \iqfirm{any} \iqlevel{junior}]\label{iq:iv3:python-and-pandas:3}
On CPython 3.10, what two lines does \cref{lst:iv3:pp:is} print? What is the difference between \texttt{is} and \texttt{==}, and when should you use \texttt{is}?
\omcode{code/interviews-3/02-python-and-pandas/python/pp_is.py}{1}{4}{Identity and equality.}\label{lst:iv3:pp:is}
\end{question}

\begin{question}[$\star$ \iqroles{QD, Dev, DS} \iqfirm{any} \iqlevel{junior}]\label{iq:iv3:python-and-pandas:4}
What does \cref{lst:iv3:pp:keys} print? Beyond ``a list can change and a tuple cannot'', what follows from that difference for dictionary keys, and for memory?
\omcode{code/interviews-3/02-python-and-pandas/python/pp_keys.py}{1}{7}{What can be a key.}\label{lst:iv3:pp:keys}
\end{question}

\begin{question}[$\star$ \iqroles{QD, Dev, DS} \iqfirm{any} \iqlevel{junior}]\label{iq:iv3:python-and-pandas:5}
What two lines does \cref{lst:iv3:pp:copy} print? A colleague wants a cheap snapshot of a book held as a dictionary of lists, taken many times a second. What do you tell them?
\omcode{code/interviews-3/02-python-and-pandas/python/pp_copy.py}{1}{9}{Two kinds of copy.}\label{lst:iv3:pp:copy}
\end{question}

\begin{question}[$\star$ \iqroles{QD, Dev, QR, DS} \iqfirm{any} \iqlevel{junior}]\label{iq:iv3:python-and-pandas:6}
What three lines does \cref{lst:iv3:pp:nan} print on CPython 3.10? What does each one tell you about taking the largest value of a price list that can hold a NaN?
\omcode{code/interviews-3/02-python-and-pandas/python/pp_nan.py}{1}{4}{NaN in plain Python.}\label{lst:iv3:pp:nan}
\end{question}

\begin{question}[$\star$ \iqroles{QD, Dev, DS} \iqfirm{any} \iqlevel{junior}]\label{iq:iv3:python-and-pandas:7}
What four lines does \cref{lst:iv3:pp:exhaust} print? A function takes an iterable of returns, computes its length with a loop, then its sum, and divides. It works on a list and gives a wrong answer on a generator, with no error. What does it return there, and what changes if the sum comes first?
\omcode{code/interviews-3/02-python-and-pandas/python/pp_exhaust.py}{2}{8}{One pass or many.}\label{lst:iv3:pp:exhaust}
\end{question}

\begin{question}[$\star$ \iqroles{QD, Dev, QR, DS} \iqfirm{any} \iqlevel{junior}]\label{iq:iv3:python-and-pandas:8}
With NumPy 2.2, what three lines does \cref{lst:iv3:pp:overflow} print? Is there any warning? Where in a data pipeline would this bite?
\omcode{code/interviews-3/02-python-and-pandas/python/pp_overflow.py}{1}{7}{Fixed-width integers.}\label{lst:iv3:pp:overflow}
\end{question}

\begin{question}[$\star$ \iqroles{QR, DS, QD, Analyst} \iqfirm{any} \iqlevel{junior}]\label{iq:iv3:python-and-pandas:9}
In \cref{lst:iv3:pp:seriesand}, with pandas 2.3, which of the lines A, B and C work, and what do they return? Explain each failure.
\omcode{code/interviews-3/02-python-and-pandas/python/pp_series_and.py}{1}{13}{Three ways to combine two conditions.\protect\hfil\protect\break}\label{lst:iv3:pp:seriesand}
\end{question}

\begin{question}[$\star\star$ \iqroles{QD, Dev, QR, DS} \iqfirm{any} \iqlevel{junior, experienced}]\label{iq:iv3:python-and-pandas:10}
On CPython 3.10, four threads each run a pure Python loop that does arithmetic, on a machine with four free cores. What speed-up do you get over one thread, and why? Does that make \texttt{count += 1} safe on a shared global? When do threads help, and what do you use when they do not?
\end{question}

\begin{question}[$\star\star$ \iqroles{QD, Dev, QR} \iqfirm{any} \iqlevel{junior}]\label{iq:iv3:python-and-pandas:11}
Write a decorator \texttt{memo} that caches a function's results by its arguments. What does it do to the number of calls of a recursive Fibonacci function for $n = 30$? What happens with a list as an argument, and are \texttt{f(1, 2)} and \texttt{f(1, b=2)} one entry or two? Why not write \texttt{cache.get(key) or f(\dots)}?
\end{question}

\begin{question}[$\star\star$ \iqroles{QD, Dev, DS} \iqfirm{any} \iqlevel{junior}]\label{iq:iv3:python-and-pandas:12}
A CSV file of 50 GB has columns \texttt{symbol}, \texttt{price\_ticks} and \texttt{qty}. On a machine with 16 GB of memory, in plain Python, compute the total notional per symbol. What does your memory use depend on? If a generator holds the file open and its caller stops reading halfway, when does the file close?
\end{question}

\begin{question}[$\star\star$ \iqroles{QD, Dev} \iqfirm{any} \iqlevel{junior}]\label{iq:iv3:python-and-pandas:13}
What two lines does \cref{lst:iv3:pp:classattr} print, and why do they look inconsistent? How would \texttt{\_\_slots\_\_} change this class, and what do you gain with a million instances?
\omcode{code/interviews-3/02-python-and-pandas/python/pp_class_attr.py}{1}{15}{An attribute of the class.}\label{lst:iv3:pp:classattr}
\end{question}

\begin{question}[$\star\star$ \iqroles{QD, Dev} \iqfirm{any} \iqlevel{junior}]\label{iq:iv3:python-and-pandas:14}
Write a context manager, as a class and then with \texttt{contextlib}, that logs ``acquire'' on entry and ``release'' on exit, even when the body raises. What happens if \texttt{\_\_exit\_\_} returns true? What goes wrong in the \texttt{contextlib} version if you leave out \texttt{try} and \texttt{finally}?
\end{question}

\begin{question}[$\star\star$ \iqroles{QD, Dev} \iqfirm{any} \iqlevel{junior}]\label{iq:iv3:python-and-pandas:15}
You write a dataclass \texttt{Order} with a field \texttt{tags: list = []}. What happens, and when? Fix it. Then you put \texttt{Order} objects in a set, and it fails too. Why? What does \texttt{frozen=True} give you, and what does it not protect?
\end{question}

\begin{question}[$\star\star$ \iqroles{QD, Dev, DS} \iqfirm{any} \iqlevel{junior}]\label{iq:iv3:python-and-pandas:16}
You have a list of trades, each a pair of a symbol and a notional. Sort them by symbol ascending and, within a symbol, by notional descending. Trades that tie must keep their original order. Show a way that works even when the second key cannot be negated, such as a string. In what order do you sort?
\end{question}

\begin{question}[$\star\star$ \iqroles{QR, DS, QD} \iqfirm{any} \iqlevel{junior}]\label{iq:iv3:python-and-pandas:17}
In \cref{lst:iv3:pp:broadcast}, \texttt{r} holds returns with one row per day and one column per stock, and the aim is to subtract each day's cross-sectional mean. With NumPy 2.2, what do the two prints give? Which is right, how can you tell at a glance, and what would have happened with three stocks in place of two?
\omcode{code/interviews-3/02-python-and-pandas/python/pp_broadcast.py}{1}{5}{Subtracting each row's mean.\protect\hfil\protect\break}\label{lst:iv3:pp:broadcast}
\end{question}

\begin{question}[$\star\star$ \iqroles{QR, DS, QD} \iqfirm{any} \iqlevel{junior}]\label{iq:iv3:python-and-pandas:18}
What does \cref{lst:iv3:pp:views} print with NumPy 2.2? Which of \texttt{b}, \texttt{c} and \texttt{d} are views? You keep \texttt{huge[:10]} from an array of 10 GB. What happens to the memory? In pandas 2.3, what does \texttt{df[df.x > 0][\textquotesingle y\textquotesingle] = 1} do to \texttt{df}?
\omcode{code/interviews-3/02-python-and-pandas/python/pp_views.py}{1}{10}{Views and copies.}\label{lst:iv3:pp:views}
\end{question}

\begin{question}[$\star\star$ \iqroles{QR, DS, QD, Analyst} \iqfirm{any} \iqlevel{junior}]\label{iq:iv3:python-and-pandas:19}
A DataFrame of trades has columns \texttt{symbol}, \texttt{day}, \texttt{price} and \texttt{qty}, with rows (A, 1, 10, 100), (A, 1, 11, 300), (B, 1, 20, 50) and (A, 2, 12, 0). Compute the volume-weighted average price per symbol and day without \texttt{apply}. What does it give for each group? What does \texttt{groupby(\dots).price.mean()} give for (A, 1), and why is that wrong?
\end{question}

\begin{question}[$\star\star$ \iqroles{QR, DS, QD, QT} \iqfirm{any} \iqlevel{junior}]\label{iq:iv3:python-and-pandas:20}
For one symbol, trades happen at integer times 1, 5 and 10, and quotes at times 2, 5 and 9. Attach to each trade the latest quote with \texttt{pd.merge\_asof}. Which quote time does each trade get? What changes with \texttt{allow\_exact\_matches=False}, and with \texttt{direction=\textquotesingle forward\textquotesingle}, and which of them is a bug in a backtest? What happens if the input is unsorted?
\end{question}

\begin{question}[$\star\star$ \iqroles{QR, DS, QD, QT} \iqfirm{hedge fund, market maker} \iqlevel{junior}]\label{iq:iv3:python-and-pandas:21}
Ticks arrive at 09:30:10 at price 10, 09:30:50 at 12, 09:31:20 at 11 and 09:33:05 at 13. Build one-minute open, high, low and close bars with \texttt{resample}. What is in the bar for 09:30, and under what time does pandas label it by default? Why does that label leak the future into a backtest, and what is the fix? What should the bar for 09:32 hold?
\end{question}

\begin{question}[$\star\star$ \iqroles{QR, DS, QD, Analyst} \iqfirm{any} \iqlevel{junior}]\label{iq:iv3:python-and-pandas:22}
What three lines does \cref{lst:iv3:pp:align} print with pandas 2.3? Say you have two stocks' prices and one stock skipped a day: which of the three is the bug that produces no error, and what does it do to a pair's returns?
\omcode{code/interviews-3/02-python-and-pandas/python/pp_align.py}{1}{7}{Index alignment.}\label{lst:iv3:pp:align}
\end{question}

\begin{question}[$\star\star$ \iqroles{QR, QT, DS} \iqfirm{hedge fund, market maker, bank} \iqlevel{junior}]\label{iq:iv3:python-and-pandas:23}
A backtest takes independent normal daily returns with zero mean, sets \texttt{signal = np.sign(ret)} and computes \texttt{pnl = signal * ret}. What annual Sharpe ratio (mean over standard deviation, times $\sqrt{252}$) do you get, for a long sample? Derive it. What is the one-word fix in pandas, and what does the Sharpe become? How would you make this kind of bug hard to write in a research library?
\end{question}

\begin{question}[$\star\star$ \iqroles{QR, DS, QD, Analyst} \iqfirm{any} \iqlevel{junior}]\label{iq:iv3:python-and-pandas:24}
What four lines does \cref{lst:iv3:pp:nanpandas} print with pandas 2.3? A colleague runs \texttt{prices.fillna(0)} before computing returns. What goes wrong, and what would you do instead?
\omcode{code/interviews-3/02-python-and-pandas/python/pp_nan_pandas.py}{1}{8}{NaN in pandas.}\label{lst:iv3:pp:nanpandas}
\end{question}

\begin{question}[$\star\star$ \iqroles{QD, DS, QR} \iqfirm{any} \iqlevel{junior, experienced}]\label{iq:iv3:python-and-pandas:25}
A table has 100 million rows: \texttt{symbol}, a string with 4\,000 distinct values, \texttt{price}, a float, and \texttt{qty}, an integer whose range depends on the data. In pandas 2.3, which types would you choose to save memory? About how much does each column take? A colleague suggests \texttt{float32} for prices and for running volume totals. Show what that breaks.
\end{question}

\begin{question}[$\star\star$ \iqroles{QD, Dev, QR} \iqfirm{any} \iqlevel{junior}]\label{iq:iv3:python-and-pandas:26}
Write pytest tests for a function that computes a volume-weighted price per group, using \texttt{parametrize}. Why does \texttt{0.1 + 0.2 == 0.3} fail, and what does \texttt{pytest.approx} allow by default? Does \texttt{1e-13 == pytest.approx(0)} pass, and does \texttt{1e-9 == pytest.approx(0)}? Which edge cases do you add?
\end{question}

\begin{question}[$\star\star$ \iqroles{QD, Dev, QR} \iqfirm{any} \iqlevel{junior}]\label{iq:iv3:python-and-pandas:27}
What four lines does \cref{lst:iv3:pp:round} print on CPython 3.10? A fee must be rounded to cents with halves going up, and the value is 2.675 exactly. How do you do it?
\omcode{code/interviews-3/02-python-and-pandas/python/pp_round.py}{1}{6}{Rounding.}\label{lst:iv3:pp:round}
\end{question}

\begin{question}[$\star\star$ \iqroles{QR, QD, LLE} \iqfirm{any} \iqlevel{junior, experienced}]\label{iq:iv3:python-and-pandas:28}
On CPython 3.10, what do \texttt{sum([1e16, 1.0, -1e16])} and \texttt{math.fsum([1e16, 1.0, -1e16])} return, and why? Write your own compensated sum, in one pass with constant memory. Does the classic Kahan method get that list right?
\end{question}

\begin{question}[$\star\star$ \iqroles{QD, Dev, LLE} \iqfirm{crypto, market maker} \iqlevel{junior, experienced}]\label{iq:iv3:python-and-pandas:29}
An \texttt{asyncio} process has a heartbeat coroutine that must send every 10 ms, and a handler that sometimes does 200 ms of blocking work inside a coroutine. What happens to the heartbeat, and why? Fix it for blocking input or output, and for pure computation. What happens if you call a coroutine and forget the \texttt{await}?
\end{question}

\begin{question}[$\star\star\star$ \iqroles{QR, DS, QD} \iqfirm{any} \iqlevel{junior, experienced}]\label{iq:iv3:python-and-pandas:30}
A table holds 10 million rows of \texttt{symbol}, \texttt{x} (a factor return) and \texttt{y} (the stock's return), for 5\,000 symbols, with NaN in either column. In pandas, without \texttt{apply} and without a loop over symbols, compute the least-squares slope of \texttt{y} on \texttt{x} per symbol. For a symbol whose rows are (1, 2), (2, 4), (3, 6) and (100, NaN), what is the right slope? What do you get if you demean each column on its own and then sum, letting pandas skip the NaN? What should a symbol with one row return?
\end{question}

\begin{question}[$\star\star\star$ \iqroles{QR, DS, QD, QT} \iqfirm{hedge fund, market maker} \iqlevel{junior, experienced}]\label{iq:iv3:python-and-pandas:31}
A quote table has \texttt{t}, \texttt{symbol}, \texttt{bid} and \texttt{ask}, across several days. Some rows have a zero side or a crossed market, some share a timestamp, some arrive out of order, and one mid of 1\,000 sits among mids of 100. In pandas, write a cleaning step and then one-minute log returns of the mid, per symbol, with no return across a night. In which order do the steps go? Why must an outlier test not use a centred window? If you reject a mid more than 5\% from the median of the 20 before it, what happens to a true jump of 6\%?
\end{question}

\begin{question}[$\star\star\star$ \iqroles{QD, QR, LLE, QT} \iqfirm{any} \iqlevel{junior, experienced}]\label{iq:iv3:python-and-pandas:32}
Values arrive one at a time. Write an exponentially weighted mean with constant memory per stream that matches pandas 2.3's \texttt{ewm(alpha=a).mean()} exactly, NaN included. For $a = 0.1$ and the values 1, 2, 3, what does that give at the third value, and what does \texttt{adjust=False} give? When do the two agree? Now ticks arrive at uneven times: change it so that a weight halves every fixed period of time. What does your version do when two ticks share a timestamp, and after a gap of a day?
\end{question}

\begin{question}[$\star\star\star$ \iqroles{QR, QD, Risk} \iqfirm{any} \iqlevel{junior, experienced}]\label{iq:iv3:python-and-pandas:33}
With NumPy, draw $n$ samples of $k$ standard normals with correlation matrix $C$. A colleague computes \texttt{L = np.linalg.cholesky(C)}, draws \texttt{z} with one sample per row, and returns \texttt{z @ L}. For two variables with correlation 0.8, what covariance matrix do the samples have? Give the correct line and prove it. What happens with pairwise correlations 0.9, 0.9 and $-0.9$, why can no data produce them, and what do you do? What happens with a correlation of exactly 1?
\end{question}

\begin{question}[$\star\star\star$ \iqroles{QR, QD, Risk} \iqfirm{bank, hedge fund, market maker} \iqlevel{junior, experienced}]\label{iq:iv3:python-and-pandas:34}
Under Black--Scholes, take spot and strike both at 100, a 5\% rate, 20\% volatility and one year: the call is worth 10.4506. Price it by Monte Carlo in NumPy, with antithetic pairs and a standard error, so that the path count is not limited by memory. What goes wrong if you treat the two members of each pair as independent when computing the error? Is the error then too large or too small, and does that depend on the payoff? Give one exact check on your code that needs no closed-form option price.
\end{question}

\begin{question}[$\star\star\star$ \iqroles{QD, QT, QR} \iqfirm{any} \iqlevel{junior, experienced}]\label{iq:iv3:python-and-pandas:35}
Signed orders arrive, and each is cut so that the position stays within $\pm L$; the next order starts from the cut position. A colleague computes the position as \texttt{np.clip(np.cumsum(orders), -L, L)}. Give the smallest example you can where that is wrong. Write it correctly. Can the two-sided version be vectorised? If only an upper cap applies, is there a closed form with \texttt{cumsum}?
\end{question}

\begin{question}[$\star\star\star$ \iqroles{QD, Dev, LLE} \iqfirm{any} \iqlevel{junior, experienced}]\label{iq:iv3:python-and-pandas:36}
On CPython 3.10, a service builds a new \texttt{Pricer} per request, as in \cref{lst:iv3:pp:growth}. Memory grows and never falls, and \texttt{gc.collect()} frees nothing. What holds the objects? After 50 requests, what does \texttt{Pricer.price.cache\_info()} report, and why is the cache useless as well as harmful? How do you find this kind of growth in a live process, and what can your tool not see? Fix it. Is \texttt{maxsize=128} a fix?
\omcode{code/interviews-3/02-python-and-pandas/python/iv_growth.py}{2}{15}{A cache on a method.}\label{lst:iv3:pp:growth}
\end{question}

\begin{question}[$\star\star\star$ \iqroles{QR, QD, DS} \iqfirm{any} \iqlevel{junior, experienced}]\label{iq:iv3:python-and-pandas:37}
A backtest sets its random seed and reads the same data, yet two runs print different numbers. A reduced form is \cref{lst:iv3:pp:hashseed}. On CPython 3.10, why can two runs of this program print different lines, in both print statements? Would the same happen with integer symbol codes? How do you fix it, and why is setting \texttt{PYTHONHASHSEED} the wrong fix? Which other sources of this problem would you check?
\omcode{code/interviews-3/02-python-and-pandas/python/pp_hashseed.py}{1}{5}{Two runs, two answers.}\label{lst:iv3:pp:hashseed}
\end{question}

\begin{question}[$\star\star\star\star$ \iqroles{QR, DS, QD} \iqfirm{hedge fund} \iqlevel{experienced}]\label{iq:iv3:python-and-pandas:38}
A table has \texttt{date}, \texttt{symbol}, \texttt{price} and \texttt{signal}, for 3\,000 stocks over 2\,500 days, with NaN in \texttt{signal}. For each date, compute the Spearman rank correlation between the signal and the next day's return of each stock, without \texttt{apply} and without a loop over dates. Why must the next return not come from \texttt{df.price.shift(-1)} on the sorted table? If the signal and the next return are jointly normal with correlation 0.1, what average do you expect? How noisy is one day's value with 500 stocks? What changes in the t-statistic if the target becomes a five-day return?
\end{question}

\begin{question}[$\star\star\star\star$ \iqroles{QD, LLE, DS} \iqfirm{any} \iqlevel{experienced}]\label{iq:iv3:python-and-pandas:39}
With NumPy only, write an as-of join. You get arrays \texttt{trade\_sym}, \texttt{trade\_t}, \texttt{quote\_sym} and \texttt{quote\_t}: symbols are integer codes from 0 to 4\,999, times are \texttt{int64} nanoseconds since 1970, near $1.8 \times 10^{18}$, and nothing is sorted. For each trade, return the row of the latest quote of the same symbol at or before its time, or $-1$. Quotes with equal time and symbol: the last one to arrive wins. There may be no loop over symbols. Why does the key \texttt{sym * 10**19 + t}, or any multiplier large enough for the times, fail? What does your version cost?
\end{question}

\begin{question}[$\star\star\star\star$ \iqroles{QR, QD, LLE} \iqfirm{any} \iqlevel{experienced}]\label{iq:iv3:python-and-pandas:40}
\texttt{x} and \texttt{y} are NumPy arrays of daily returns of about 1\%, with NaN. For each row, compute the slope of \texttt{y} on \texttt{x} over the last $w$ rows, using only rows where both are present, or NaN if fewer than \texttt{min\_obs} rows qualify. Use no Python loop over rows, and $O(n)$ time. The usual method differences a cumulative sum. A bad tick then puts $10^9$ in \texttt{x} at row 10 of 1\,000. With $w = 50$, what does that method give at row 999, where the true slope is near 2? Can you avoid the damage and stay $O(n)$? What does pandas 2.3.3's \texttt{rolling(50).var()} show at that row?
\end{question}
